\section{The center-focus problem}

The plane differential system
\begin{align}\label{system}
\dot{x} &= P(x,y), \quad
\dot{y}= Q(x,y)
\end{align}
is said to have a \emph{center} at the singular point $(0,0)$, if in a sufficiently small neighbourhood of this point all orbits are closed. Consider the scalar differential equation
\begin{align}\label{ode}
\frac{\dy}{\dx} + f(x,y) = 0,\quad x\in [0,1]
\end{align}
in which $f(x,0) =0, \forall x\in [0,1]$. The equation~\eqref{ode} is said to have a \emph{center} at $y=0$, if all solutions $y(x)$ starting near the origin, satisfy $y(0)=y(1)$ (the interval $[0,1]$ can be replaced by any closed interval).


\emph{Note on the terminology.} We do not specify here the category to which belong $P, Q, f$. They will be either analytic or polynomial, depending on the context. The base field will be either $\R$ or $\Cbb$ depending on the context too. Most results will be valid for both. Thus, the definition of a center for~\eqref{ode} is the same in the real and in the complex case. In the case of an analytic complex plane vector field~\eqref{system} the ``complex'' definition of a center is less straightforward. We say that the origin is a non-degenerate center, if the vector field has an analytic first integral with a Morse critical point at the origin. If this is the case, we shall also say that~\eqref{system} has a Morse singular point, e.g.~\cite{cerlin96,dul08}. We recall therefore
\begin{defi}
The analytic complex vector field~\eqref{system} is said to have a Morse singular point, if it allows an analytic first integral in a neighbourhood of this point, which has a Morse type singularity.
\end{defi}

If~\eqref{system} has a Morse singular point, then the linear part of~\eqref{system} is diagonalisable with non-zero eigenvalues, that is to say the singular point of the vector field is non-degenerate.

An example is the saddle $x'=x, y'=-y$ which has an analytic first integral $xy$ of Morse type, and hence a Morse critical point. Of course, it is linearly equivalent (over $\Cbb$) to $x'=y,y'=-x$ with first integral $x^2+y^2$ which is the usual linear real center. The advantage to study Morse critical points over $\Cbb$ is that we can use complex analysis and complex algebraic geometry. This is the point of view adopted in these notes.

The two equations~\eqref{system} and~\eqref{ode} are closely related. First, a polar change of variables transforms a plane system~\eqref{system} with a center, to equivalent equation of the form~\eqref{ode} with a center along the interval $[0,2\pi]$. Second, if the family of functions $f({\,\cdot\,},y), x\in [0,1]$ is replaced by its Fourier series $\hat{f}({\,\cdot\,},y)$ (so $\hat{f}(x+1,y) = \hat{f}(x,y)$) and the equation~\eqref{ode} has a center at $y=0$, then the new system
\begin{align}\label{ode1}
\frac{\dy}{\dx} + \hat{f}(x,y) = 0, \quad(x,y)\in \R / \Z \times \R
\end{align}
will have all its orbits starting near the periodic solution $y=0$ on the cylinder $ \R / \Z \times \R$, periodic too. Of course, if the smooth function $f$ is non-periodic, then the function $\hat{f}$ is only piece-wise continuous in $x$. The transport map of~\eqref{ode} along $[0,1]$ becomes a return map for~\eqref{ode1} and the definition of a limit cycle for~\eqref{ode} is straightforward too. Actually, the scalar equation~\eqref{ode} in which $f$ is a regular function, should be considered as a simplified model of the eventually singular equation
\[
\frac{\dy}{\dx} = \frac{P(x,y)}{Q(x,y)}.
\]

We resume the above considerations in the following definitions, which make sense both on $\R$ or $\Cbb$:
\begin{defi}\label{def1}
Let $ \varphi = \varphi({\,\cdot\,}; x_0,y_0)$ be the general solution of the equation $\dy + f(x,y)\dx=0$ with initial condition $y_0=\varphi(x_0; x_0,y_0)$, on the interval $ [x_0,x_1]$.
%%A: desactivation de l'environnement, description
%%\begin{description}
\begin{enumerate}\romanenumi
\item \label{list1.2.1}The solution $ \varphi = \varphi({\,\cdot\,}; x_0,y_0)$ is said to be periodic iff $\varphi(x_1; x_0,y_0) = y_0$
\item \label{list1.2.2}The solution $ \varphi = \varphi({\,\cdot\,}; x_0,y_0)$ is said to be a limit cycle, provided that it is periodic and isolated, that is to say there is a neighbourhood of its orbit on $S^1\times \R$ free of periodic solutions.
\item \label{list1.2.3}the map $y \mapsto \varphi(x_1; x_0,y) $ is the first return map of~\eqref{ode} in a neighbourhood of $ y=y_0$.
\item \label{list1.2.4}The equation~\eqref{ode} defines a center in a neighbourhood of the periodic solution $\varphi$ provided that the first return map is the identity map in a neighbourhood of $y_0$. If the return map is not the identity map, then we say that~\eqref{ode} defines a focus at the periodic solution $\varphi$.
\end{enumerate}
%%\end{description}
\end{defi}

The center focus-problem for the equation~\eqref{ode} or~\eqref{system} is, roughly speaking, to distinguish between a center and a focus. The algebro-geometric content of the problem is as follows. Suppose, that~\eqref{ode} is polynomial, more precisely
\begin{align}\label{ode2}
\frac{\dy}{\dx} + \sum_{i=0}^m a_i(x) y^{i+1},\quad a_i \in \Cbb[x], \deg a_i \leq n, x \in [0,1].
\end{align}

The first return map $y \mapsto \varphi(1;0,y)$ is well defined and analytic near the periodic solution $y=0$, and moreover
\[
\varphi(1;0,y) = y + \sum_{n=1}^\infty c_n(a) y^{n+1}.
\]
As we shall see in Theorem~\ref{brf}, under the condition $a_1=0$, the coefficients\linebreak $c_n=c_n(a), n\geq 1$, are \emph{polynomials} in the coefficients of $a_j = a_j(x)$, $j\leq n$. The condition that $\varphi(1;0,{\cdot\,})$ is the identity map determines an infinite number of polynomial relations $\{ c_n(a)=0 \}$ on the coefficients of the polynomials $a_j$. By the Hilbert basis theorem, only a finite number of them are relevant, and they define an algebraic variety (the so called center variety $\mathcal C_{m,n}$) in the vector space of all coefficients of the polynomials $a_j$. The problem is therefore (as formulated by Lins Neto~\cite{lins14} in the context of a polynomial foliation induced by~\eqref{system}):
\[
\text{\emph{Describe the irreducible components of $\mathcal C_{m,n}$.}}
\]

%%A: je rentre une subsection* pour l'intitule de type center, faisant l'introduction du plan de l'article:
%%\begin{center}
%\emph{Describe the irreducible components of} $\mathcal C_{m,n}$.
%\end{center}
%%~ N: Non, il s'agissait de la fin de la phrase précédente
%\pagebreak
%\subsection*{Describe the irreducible components of \texorpdfstring{$\mathcal C_{m,n}$}{Cm,n}}
%%\vspace{1cm}
\goodbreak
The content of the paper is as follows.

In Section~\ref{section2} we give first an explicit formula for the general solution of the equation
\[
\frac{\dy}{\dx} + \sum_{i=0}^n a_i(x) y^{i+1} = 0,\quad 0\leq x \leq 1
\]
in terms of iterated path integrals, see Theorem~\ref{firstintegral}. As a by-product we obtain a formula for the first return map, and explicit center conditions found first by Brudnyi, see Theorem~\ref{brf}.



Section~\ref{section3} is devoted to the perturbation theory of the integrable Abel equation
\[
\frac{\dy}{\dx} = a(x) y^2
\]
with first integral
\[
H= \frac1y + A(x),\quad A(x)= \int a(x) \dx.
\]
It is assumed that $A(0)=A(1)$, so the equation has a center along $[0,1]$. We are interested in the number of limit cycles (isolated solutions, such that $y(0)=y(1)$) which the perturbed equation $
\frac{\dy}{\dx} = a(x) y^2 + \dots $ can have. The center-focus problem for this perturbed equation leads to a well known \emph{ polynomial moment problem}. Under general assumptions this problem has an elegant solution, due to Colin Christopher, Theorem~\ref{ccth}, which is presented here in the setting of Abelian integrals of dimension zero.

In Section~\ref{section4} we study irreducible components of the center variety of the polynomial Abel equation (on the interval $[0,1]$)
\begin{align}\label{abel1}
\frac{\dy}{\dx} = p(x) y^2 + q(x) y^3
\end{align}
as well center variety of the related Liénard equation (by ``center'' we mean the usual Morse center in a neighbourhood of the origin in $\mathbb C^2$)
\begin{align}\label{lienard1}
\dot{x} = y, \quad \dot{y} = -q(x) - y p(x).
\end{align}
In Section~\ref{section41} we prove that the set of Abel equations coming from ``pull back'' provide irreducible components of the center set, Theorem~\ref{pullback}. These results are inspired by previous contributions of Movasati.

In Sections~\ref{section42} we revisit the classical center-focus for quadratic vector fields, with special attention to the $Q_4$ component of the center set.

In Section~\ref{section4b} we give a full description of the center set of Liénard type equations~\eqref{lienard1}. These results belong mainly to Cherkas and Christopher, but we present them in the broader context of the present notes. In particular, the base field will be $\Cbb$, see Theorem~\ref{th2a} and~\ref{th2}. The centers found in this way are always of ``pull back'' type. This suggests that the only centers of the related Abel equation~\eqref{abel1} are of ``pull back'' type too, which is the content of the so called Composition Conjecture for the Abel equation~\eqref{abel1}~\cite[p.~444]{bry10} to be discussed in Section~\ref{section43}. In this last section we show, however, that there are scalar Abel equations with a center along $[0,1]$, which can not be obtained by a ``pull back''. These equations have a Darboux type first integral, and their construction is inspired by the study of the $Q_4$ component in Section~\ref{section42}. Among them we find the recent counter-example to the Composition Conjecture mentioned above, found first by Giné, Grau and Santallusia~\cite{ggs18}.\footnote{\label{myfootnote}
The present paper is an extended version of two lectures given during the Zagreb Dynamical Systems Workshop, October 22-26, 2018.}%un espace en trop était introduit par le retour à la ligne.

\section*{Acknowledgement}
The author is obliged to Jean-Pierre Françoise for the illuminating discussions on the center-focus problem. I thank also the anonymous referees for several suggestions, which helped to improve the text.

\section{The first return map and the Brudnyi formula}\label{section2}
In this section we shall describe the return map of~\eqref{ode2} as a power series involving iterated path integrals. We prove an explicit formula, due to Brudnyi~\cite{brud06}, which amounts to solve the differential equation. The classical approach to do this is by the Picard iteration method. If $y_0$ is the initial condition at $x_0$ of the differential equation
\[
\dy = f(x,y) \dx
\]
then the Picard iteration is
\[
y_{n+1}(x) = y_0 + \int_{x_0}^x y_n(t) \dt
\]
where $y_n$ tends to the solution of the equation as $n\to \infty$. We illustrate this on the example $\dy = y \dx$. If $y_0$ is the initial condition at $x = 0$ then
\begin{align*}\label{exp}
y_1(x) & = y_0 + \int_0^x y_0 \dt\\
y_2(x) & = y_0 + \int_0^x y_1(t) \dt = y_0 + \int_0^x y_0 \dt + \iint_{0\leq t_2\leq t_1\leq x} y_0 \dt_1 \dt_2.
\end{align*}
As
\[
\int \dots \int_{0 \leq t_n\leq \dots \leq t_1\leq x}y_0 \dt_1 \dots \dt_n = y_0 \frac{x^n}{n!}
\]
we get $y(x) = y_0 e^x$ as expected. The multiple (or iterated) integrals above appear in a similar way in the non-autonomous linear $\dy=a(x) y \dx$, or even non-linear case $\dy=f(x,y) \dx$. The non-linear case is more involved, it is reduced to the linear one, but after introducing infinitely many new variables $y, y^2, y^3, \dots $. To get around this reduction we shall use a simple Ansatz, for which we need a formal definition of iterated integral.



Let $Ass_\omega$ be the graded free associative algebra generated by the infinite dimensional vector space of differential one-forms $\omega = a(x,y) \dx$, $a \in \Cbb\{x,y\}$. Its elements are non-commutative polynomials in such one-forms. The differential operator
\begin{gather*}
D : Ass^1_\omega \to Ass^1_\omega
\\
D(a(x,y) \dx) = \frac{\partial}{\partial y} a(x,y) \dx
\end{gather*}
induces a differential operator on $ Ass_\omega$ which acts by the Leibnitz rule. The readers familiar with the Picard--Lefschetz theory will recognize in $D$ an avatar of the covariant derivative of an Abelian differential on the level sets $\{y=c \}_c$.

To save brackets, it is convenient to introduce the following notation
\begin{equation}\label{Domega}
D\omega_1 \omega_2 \dots \omega_n = D(\omega_1 \omega_2 \dots \omega_n)
\end{equation}
so that (using brackets)
\begin{align*}
D \omega_1 \omega_2 =D (\omega_1 \omega_2) &= (D \omega_1) \omega_2 + \omega_1(D \omega_2).
\\
\intertext{and}
D \omega_1 D \omega_2 = D (\omega_1 D \omega_2) &= (D \omega_1)(D \omega_2) + \omega_1 (D^2 \omega_2).
\end{align*}

If we use the notation
\[
D^k\omega = \omega ^{(k)}
\]
then
\[
D \omega_1 \omega_2 = \omega_1' \omega_2+ \omega_1 \omega_2'
\]
and
\[
D \omega_1 D \omega_2 = (\omega_1 \omega_2')' = \omega_1' \omega_2' + \omega_1 \omega_2''.
\]
For $\omega_1\omega_2 \dots \omega_n \in Ass_\omega^n$, $\omega_k= \varphi_k(x,y) \dx$, define the iterated integral $\int_{x_0}^x \omega_1\omega_2 \dots \omega_n$ of length $n$, as equal to
\begin{equation}\label{multipleintegral}
\int \dots \int_{x_0 \leq t_n\leq \dots \leq t_1\leq x} \varphi_1(t_1,y) \dots \varphi_n(t_n,y) \dt_1 \dots \dt_n.
\end{equation}
The iterated integral allows also a recursive definition (hence the name):
\begin{equation}\label{iteratedintegral}
\int_{x_0}^x \omega_n\omega_{n-1} \dots \omega_1 = \int_{x_0}^x (\varphi_n (t) \int_{x_0}^t \omega_{n-1} \dots \omega_1)\dt
\end{equation}
where in the case $n=1$ we have the Riemann integral $\int_{x_0}^x \omega_{1}$. We note, that the usual notation for the multiple integral~\eqref{multipleintegral} is $\int_{x_0}^x \omega_n\omega_{n-1} \dots \omega_1$ on the place of $\int_{x_0}^x \omega_1\omega_2 \dots \omega_n$, see Chen~\cite{chen77} or Hain~\cite{hain87}. The reason to prefer the definition~\eqref{iteratedintegral} is that it is better adapted to applications in differential equation, e.g.~\cite{gavr05}. Recall in this context, that
\[
\int_{x_0}^x \omega_n\omega_{n-1} \dots \omega_1 = (-1)^n \int^{x_0}_x \omega_1\omega_2 \dots \omega_n.
\]
For a short summary of properties of iterated integrals which we use, see~\cite[Appendix]{gavr05}, \cite[Section~2]{gmn09}.


\begin{theo}\label{firstintegral}
With the notation~\eqref{Domega}, a first integral of the differential equation $\dy + f(x,y) \dx = 0$ is given by the following recursively defined convergent series
%%\boxed{}
\begin{align}\label{first2}
\varphi(x_0;x,y) = y + \int_{x_0}^x \omega + \int_{x_0}^x \omega D \omega + \int_{x_0}^x \omega D \omega D \omega +
\dots 
\end{align}
where
\[
\omega = f(x,y) \dx.
\]
The general solution of~\eqref{ode} with initial condition $(x_0,y_0)$ is given by
\[
y= \varphi(x;x_0,y_0).
\]
\end{theo}

\begin{exam}
In the linear case
\[
y' + \alpha y = 0 \iff \dy + \alpha y \dx = 0
\]
we obtain
\begin{align*}
\varphi(x_0;x,y) &= y \left(1+ \alpha \int_{x_0}^x \dx + \alpha^2 \int_{x_0}^x \dx.\dx + \dots\right)\\
&= y(1 + \alpha(x-x_0) + \alpha^2 \frac{(x-x_0)^2}{2} + \dots =\;y e^{\alpha (x-x_0) }
\end{align*}
and the general solution is
\[
y = \varphi(x;x_0,y_0) =\; y_0 e^{\alpha (x_0-x) }.
\]
In the quadratic case
\[
\dy + 2 x y^2 \dx = 0, \omega = 2 x y^2 \dx
\]
we compute recursively
\begin{align*}
\int_{x_0}^x \omega &= \int_{x_0}^x 2 x y^2 \dx = x^2-x_0^2 \\
\int_{x_0}^x \omega D\omega & = \int_{x_0}^x 2 x y^2 \dx. 4 x y \dx = y^3 (x^2-x_0^2)^2 \\
\int_{x_0}^x \omega D\omega \dots D\omega &= (x^2-x_0^2)^n.
\end{align*}
Therefore we get the first integral
\[
\varphi(x_0;x,y) = y + y^2 (x^2-x_0^2) + y^3 (x^2-x_0^2)^2 + \dots
\]
and the corresponding general solution is
\begin{align*}
y= \varphi(x ; x_0,y_0) & = y_0 + y_0^2 (x_0^2-x^2) + y^3 (x^2_0-x^2)^2 + \dots \\
& = \frac{y_0}{1-y_0(x_0^2-x^2)}.
\end{align*}
\end{exam}

\begin{proof}[Proof of Theorem~\ref{firstintegral}.]
We first verify, that for every fixed $x_0$, the function $\varphi(x_0;x,y)$ is a first integral :
\begin{align*}
d\varphi(x_0;x,y)&=\frac{ \partial}{\partial x} \varphi(x_0;x,y) \dx + \frac{ \partial}{\partial y} \varphi(x_0;x,y) \dy \\
&= \omega + \omega \int_{x_0}^x D \omega + \omega \int_{x_0}^x D \omega D \omega + \omega \int_{x_0}^x D \omega D \omega D \omega + \dots \\
&\qquad\qquad+ \left(1 + \int_{x_0}^x D \omega + \int_{x_0}^x D \omega D \omega + \int_{x_0}^x D \omega D \omega D \omega + \dots\right) \dy \\
&= (\omega + \dy) \frac{ \partial}{\partial y} \varphi(x_0;x,y) = 0.
\end{align*}

As $\varphi(x_0;x_0,y_0) = y_0$ then the level set $\{ (x,y) : \varphi(x_0;x,y) = y_0 \}$ contains both $(x_0,y_0)$ and $(x,y)$. By symmetry
\[
y= \varphi(x;x_0,y_0)
\]
is the solution of~\eqref{ode} with initial condition $y(x_0)=y_0$. The convergency proof is by standard a priori estimates (omitted)
\end{proof}

Note that for fixed $x_0,x_1$ the two return maps
\[
y \mapsto \varphi(x_1;x_0,y), \quad y \mapsto \varphi(x_0;x_1,y)
\]
are mutually inverse. Therefore $\varphi(x_1;x_0,{\cdot\,}) = \id$ if and only if $ \varphi(x_0;x_1,{\cdot\,}) = \id$. Using Theorem~\ref{firstintegral} we can give explicit center conditions. Assume that
\[
f(x,y) = \sum_{i=1}^\infty a_i(x) y^{i+1}
\]
and develop the return map $ \varphi(x_0;x_1,y) $ as a power series in $y$
\begin{equation}\label{returnmap}
\varphi(x_0;x_1,y) = y + \sum_{n=1}^\infty c_n(a) y^{n+1}.
\end{equation}
If we denote, by abuse of notations, $a_{i}= a_{i}(x) \dx $ then we get for the first few coefficients $c_n(a)$ (compare to~\cite[p.~450]{bry10})
\begin{align*}
c_1(a) & = \int_{x_0}^{x_1} a_1 \\
c_2(a) & = \int_{x_0}^{x_1} a_2 + 2 a_1a_1\\
c_3(a) & = \int_{x_0}^{x_1} a_3 + 2 a_2 a_1 + 3 a_1a_2 + 6 a_1^3 \\
c_4(a) & = \int_{x_0}^{x_1} a_4 + 2 a_3 a_1 + 3 a_2^2 + 4 a_1a_3 + 6a_2a_1^2 + 8 a_1a_2a_1 + 12 a_1^2a_2 + 24 a_1^4
\end{align*}
and so on. The general form of the coefficients $c_n(a)$ is found immediately from Theorem~\ref{firstintegral}. We resume this in the following



\begin{theo}[Brudnyi's formula~\cite{brud06}]\label{brf}
The coefficients $c_n(a)$ of the first return map~\eqref{returnmap} for the differential equation
\[
\frac{\dy}{\dx} + \sum_{i=1}^\infty a_i(x) y^{i+1} = 0, x\in [x_0,x_1]
\]
are given by the formulae
\[
c_n(a) = \sum_{i_1+\dots + i_k=n} c_{i_1,\dots,i_k} \int_{x_0}^{x_1} a_{i_1} \cdots a_{i_k}
\]
where
\begin{align*}
c_{i_1} & = 1\\
c_{i_1,i_2} & = i_2+1 \\
c_{i_1,i_2,i_3} & = (i_3+1)(i_2+1)\\
& \;\:\vdots \\
c_{i_1,\dots,i_k} & = (i_k+1)(i_k+i_{k-1} +1) \dots (i_k + \cdots i_2 +1).
\end{align*}
\end{theo}

The above formula was deduced first by Brudnyi~\cite[p.~422]{brud06} under equivalent form, see also~\cite[Proposition~2.4]{bry10} in the case~\eqref{abel}.

\begin{coro}
The equation~\eqref{ode} has a center on the interval $[x_0,x_1]$ if and only if $c_n(a)=0$, for every integer $ n \geq 1$.
\end{coro}

\begin{exam}
Suppose that the equation
\[
\frac{\dy}{\dx} + a_1(x)y^2 + a_2(x)y^3 + \dots = 0
\]
has a center on the interval $[x_0,x_1]$. Then, using as above the notation $a_{i}= a_{i}(x) \dx $ we have
\begin{align*}
c_1&= \int_{x_0}^{x_1} a_1 = 0 \\
c_2& = \int_{x_0}^{x_1} a_2 + 2 \int_{x_0}^{x_1} a_1^2 = 0.
\end{align*}
The identity
\[
2 \int_{x_0}^{x_1} a_1^2 = \left(\int_{x_0}^{x_1} a_1\right)^2
\]
implies then, that $ \int_{x_0}^{x_1} a_2 = 0$.
If we consider more specifically the Abel equation
\begin{equation}\label{abelex}
\frac{\dy}{\dx} + a_1(x)y^2 + a_2(x)y^3 = 0
\end{equation}
then taking into consideration that $\int_{x_0}^{x_1} a_1^3 = 0$ and
\[
\int_{x_0}^{x_1} (a_2 a_1 + 3a_1a_2) = \int_{x_0}^{x_1} a_1 a_2 = 0
\]
we obtain $c_3= \int_{x_0}^{x_1} a_1 a_2 $. Therefore a necessary condition for the Abel equation~\eqref{abelex} to have a center on $[x_0,x_1]$ is
\begin{equation}\label{nc}
\int_{x_0}^{x_1} a_1 = 0,\quad \int_{x_0}^{x_1} a_2 = 0, \quad \int_{x_0}^{x_1} a_1 a_2 = 0
\end{equation}
If we suppose that $a_1$, $a_2$ are polynomials of degree at most two, these conditions are also sufficient. The case $\deg a_1, a_2 = 3$ can be studied similarly, see~\cite{bfy98,bfy99,bfy00}.
\end{exam}

In general, an obvious sufficient condition to have a center is therefore
\begin{equation}\label{ucenter}
\int_{x_0}^{x_1} a_{i_1} \cdots a_{i_k} = 0,\quad \forall i_j, k \geq 1.
\end{equation}
Centers with the property~\eqref{ucenter} were called \emph{universal} in~\cite{brud06}.

Consider, more specifically, the following equation with polynomial coefficients $a_i$
\begin{equation}\label{an}
\dy + \sum_{i=1}^n y^{i+1} a_i(x) \dx = 0,\quad a_i(x) \in K[x].
\end{equation}

\begin{theo}[{\cite[Corollary~1.20]{brud06}}]\label{universal}
The polynomial equation~\eqref{an} has an universal center on the interval $[x_0,x_1]$, if and only if, it is a pull back of some polynomial equation
\begin{equation}\label{bn}
\dy = \left(\sum_{i=1}^n b_i(\xi) y^{i+1}\right) \dd\xi,\quad b_i(\xi) \in K[\xi].
\end{equation}
via a suitable polynomial map $\xi = \xi (x)$ having the property $\xi(x_0)=\xi(x_1)$.
\end{theo}

Not all centers of~\eqref{an} are universal, as discovered recently in~\cite{ggs18}.


\section{Bifurcation functions related to Abel equation and a Theorem of Christopher}\label{section3}

In this section we study the following perturbed Abel differential equation on the interval $[0,1]$
\[
y' = a(x) y^2 - \sum_{j= 1}^\infty \varepsilon^j\left(y^2 p_j(x)+ y^3 q_j(x)\right)
\]
or equivalently
\begin{equation}\label{perturbed1}
\frac{\dy}{y^2} = a(x) \dx - \varepsilon \omega_1 - \varepsilon^2\omega_2 - \dots
\end{equation}
where
\begin{gather*}
\omega_j = (p_j(x)+ y q_j(x)) \dx,\quad
a=a(x),\quad p_j =p_j(x),\quad q_j=q_j(x)
\end{gather*}
are polynomials of degree
\[
\deg a = n, \deg p_j \leq n, \deg q_j \leq n
\]
and $\varepsilon$ is a small parameter, see~\cite{bfy98,bfy99,bfy00,chri00,yomd03}. For\linebreak $\varepsilon = 0$~\eqref{perturbed1} has a first integral
\[
H(x,y)= \frac1y + A(x), \quad A(x) = \int a(x) \dx.
\]
%%A: je rentre un env de type enonce pour l'intitule de type center, cite ci dessous:
%%\begin{center}
%%\emph{How many limit cycles has the perturbed system~\eqref{perturbed1} on the interval $[0,1]$? }
%%\end{center}

\begin{enonce*}[plain]{Question}
How many limit cycles has the perturbed system~\eqref{perturbed1} on the interval $[0,1]$?
\end{enonce*}
%%\vspace{0,5cm}

Recall from the preceding section that a solution $y(x)$ such that $y(0)=y(1)$, is called periodic on $[0,1]$. A limit cycle of~\eqref{perturbed1} on $[0,1]$ is therefore an isolated periodic solution on $[0,1]$.

The number of the limit cycles in a compact set are bounded by the number of the zeros of the so called \emph{bifurcation function}, which we define bellow. A limit cycle which remains bounded when $\varepsilon \to 0$, tends to a periodic solution of the non perturbed system. If the non-perturbed system ($\varepsilon = 0$) has a periodic solution, then necessarily $A(0)=A(1)$, which already implies that it has a center. For this reason we assume from now on that $A(0)=A(1)=0$, so that
\begin{equation*}
\dy = a(x) y^2 \dx \Longleftrightarrow \dH = 0
\end{equation*}
has a center along $0\leq x \leq 1$. The perturbed equation can be written
\begin{equation}\label{perturbed}
\dH - \varepsilon \omega_1 - \varepsilon^2\omega_2 - \dots = 0.
\end{equation}
For a solution $y(x)$, let $\mathcal P_\varepsilon$ be the first return map which sends the initial condition $y_0=y(0)$ to $y_1=y(1)$. We parameterise $\mathcal P_\varepsilon$ by $h=\frac1y = H(0,y)=H(1,y)$ and note that $\mathcal P_\varepsilon$ is analytic both in $h$ and $\varepsilon$ (close to zero). We have therefore for the first return map
%retour à la ligne en trop qui provoque un espace vertical en trop (cf article précédent)
\begin{equation}\label{return1}
\mathcal P_\epsilon (h)- h= \varepsilon^k M_k(h) + O(\varepsilon^{k+1}), \quad M_k \neq 0.
\end{equation}

The function $M_k$ is the\emph{ bifurcation function}, associated to the equation~\eqref{perturbed1}. It is also known as ``first non-zero Melnikov function''. The reader may compare this to~\eqref{first2} which is another representation of the first return map, defined for small $y$. As we shall see, the bifurcation function is globally defined. Therefore for every compact set $K$, $[0,1] \subset K \subset \R^2$ and all sufficiently small $|\varepsilon|$, the number of the limit cycles of~\eqref{perturbed1} in $K$ is bounded by the number of the zeros of the bifurcation function $M_k$ (counted with multiplicity).

$M_k$ allows an integral representation
\[
M_k(h) = \int_{\{H=h\}} \Omega_k
\]
where the integration is along the level set
\[
\{H=h\} = \{(x,y) : 1/y + A(x)= h, 0\leq x \leq 1 \}.
\]
The differential form $\Omega_k$ is computed by the classical Françoise's recursion formula~\cite{fran96,ilie98a,rous98} as follows:
\begin{quotation}
If $k=1$ then $\Omega_1 = \omega_1$, otherwise
\begin{equation}
\Omega_m = \omega_m + \sum_{i+j=m} r_i \omega_j, \quad 2 \leq m \leq k
\end{equation}
and the functions $r_i$, $1\leq i \leq k-1$ are determined successively from the identities $\Omega_i = \dd R_i + r_i \dH$.
\end{quotation}
Note that neither $r_i$ nor $R_i$ are uniquely defined. The integrals $M_i(h)$ are, however, defined unambiguously.

The first order Melnikov function $M_1$ was computed first by Lins Neto~\cite[Section~3]{lins80}, see also~\cite{bfy00}. We have
\begin{align*}
M_1(h) &= \int_{\{H=h\}} \omega_1 \\ &= \int_{\{H=h\}} p_1(x) \dx + yq_1(x) \dx\\
& = \int_0^1 p_1(x) \dx + \int_0^1 \frac{q_1(x)}{h - A(x)} \dx \\
& = \int_0^1 p_1(x) \dx + \sum_{k=0}^\infty h^{-k-1} \int_0^1 q_1(x) A^k(x)\dx.
\end{align*}
$M_1$ vanishes identically if and only if $\int_0^1 p_1(x) \dx = 0$ and
%%question ici de laisser ou non boxed ?Je la laisse ci dessous
%%\[
%%\boxed{
%%\int_0^1 q_1(x) A^k(x) \dx = 0, \; k= 0,1,2, \dots}
%%\]
\[
\int_0^1 q_1(x) A^k(x) \dx = 0, \quad k= 0,1,2, \dots
\]
which is the content of the famous polynomial moment problem for $q_1$ and $A$, solved in full generality by Pakovich and Muzychuk~\cite{pamu09}, see also~\cite{bfy98,bfy99,bfy00,chri00,yomd03}. If $M_1=0$ by the above formula we get
\[
M_2(h)= \int_{\{H=h\}} r_1 \omega_1 + \int_{\{H=h\}} \omega_2
\]
where $r_1$ is computed from the identity $\omega_1= \dR_1 + r_1 \dH$. As $\domega_1 = \dr_1 \wedge \dH$ then $\dr_1 = \omega_1'= \frac{\domega_1}{\dH}$ is the Gelfand--Leray form of $\omega_1$. From the identity $H(x,y(x,h)) \equiv h$ we have $ \frac{\partial y}{\partial h} = - y^2$ and hence
\[
r(x,y)= \int_0^x \omega_1' = - \int_0^x y^2 q(x) \dx.
\]
We conclude


\begin{prop}[{\cite[Formula~(2.8)]{gavr05}}]\label{prop3.1}
Under the hypothesis $M_1=0$ the second Melnikov function is given by the following iterated integral of length two
\begin{equation}
M_2(h) = \int_{\{H=h\}} \omega_1 \omega_1' + \int_{\{H=h\}} \omega_2
\end{equation}
where
\[
\omega_1= p_1(x) \dx + yq_1(x) \dx,\quad \omega_1'= -y^2 q_1(x) \dx,\quad \omega_2= p_2(x) \dx + yq_2(x) \dx.
\]
\end{prop}

The hypothesis $M_1=0$ is of interest for us, as it will allow to compute the tangent space to the center set at the point $(a,0)$, see the next Section~\ref{section4}. For our purposes the polynomial $a(x)$ can be taken in a general position, in which case the polynomial moment problem for $q(x),A(x)$ has the following elegant solution


\begin{theo}[\cite{chri00}]\label{ccth}
Assume that $A,q$ are complex univariate polynomials, such that
\[
A(0)=A(1)=0,\quad A'(0) \neq 0,\quad A'(1) \neq 0.
\]
The multivalued transcendental function
\begin{equation}\label{abelian}
I(h) = \int_0^1 \frac{q(x)}{h - A(x)} \dx
\end{equation}
vanishes identically, if and only if the polynomials $Q = \int q$ and $A$ satisfy the following ``Polynomial Composition Condition'' (PCC):
%% %%A: je rentre un env de type quotation pour l'intitule de type center, cite ci dessous:
%%\begin{center}
%%There exist polynomials $\tilde Q, \tilde A, W$, such that
%%\[
%A= \tilde A \circ W, Q= \tilde Q \circ W, W(0)=W(1).
%%\]
%%\end{center}
\begin{quotation}
\emph{There exist polynomials $\tilde Q, \tilde A, W$, such that}
\[
A= \tilde A \circ W,\quad Q= \tilde Q \circ W,\quad W(0)=W(1).
\]
\end{quotation}
\end{theo}

Before recalling the proof of Christopher, we put $I$ in the broader context of the Picard--Lefschetz theory.

The function $I(h)$ is well defined for sufficiently big $h$, and has an analytic continuation in a complex domain to certain multivalued function. It is in fact an Abelian integral depending on a parameter. More precisely, consider the genus zero affine~curve
\[
\Gamma_h= \left\{(x,y)\in \Cbb^2: \frac 1y +A(x) = h \right\}, \quad A(0)=A(1)=0.
\]
It is a Riemann sphere with $n+2$ removed points, provided that $h \neq 0$. The removed points correspond to $(x=x_i(h), y=\infty)$, where $A(x_i(h))\equiv h$, and to $(x=\infty,y=0)$. Given a divisor $m=P_0+P_1$ on $\Gamma_h$, where
\[
P_0 = \left(0,\frac1h\right), \quad P_1=\left(1,\frac1h\right)
\]
we define a singularized algebraic curve $\Gamma_h^{sing}$, see Figure~\ref{fig2}.

\begin{figure}
\includegraphics[scale=0.5]{figures/figure2.jpg} 
%%\includegraphics[width=8cm]{figure2.jpg}
\caption{The singularized algebraic curve $\Gamma_h^{sing}$.}\label{fig2}
\end{figure}%%\ \\*[-1.5em]



As a topological space it is just the curve $\Gamma_h$ with the two points $P_0$ and $P_1$ identified to a point $m$. The structural sheaf of $\Gamma_h^{sing}$ is the same as the structural sheaf of $\Gamma_h$, except at the point $m\in \Gamma_h^{sing}$. A function $f$ on $ \Gamma_h^{sing}$ is said to be regular, if it is regular on $\Gamma_h$, and moreover $f(P_0)=f(P_1)$. The path $[0,1]$ connecting the points $x=0$ and $x=1$ closes on the singular algebraic curve $\Gamma_h^{sing}$. The function $I(h)$ defined in~\eqref{abelian} is an Abelian integral on $\Gamma_h^{sing}$
\[
I(h) = \int_{\delta(h)} y q(x) \dx, \quad y = \frac{1}{h-A(x)}
\]
where $\delta(h) \in H_1(\Gamma_h^{sing},\Z)$ is represented by the closed loop on $\Gamma_h^{sing}$ corresponding to the interval $[0,1]$.

We note that given an arbitrary effective divisor $m=P_0+P_1+ \dots P_k$ on $\Gamma_h$, one constructs in a similar way a singularized curve $\Gamma_h^{sing}$, which is the natural framework of the generalized center problem for the Abel equation, see~\cite[Conjecture~1.7]{bfy99} and~\cite{bry10}.



\begin{proof}[Proof of Theorem~\ref{ccth}]
The homology group $H_1(\Gamma_h^{sing}, \Z)$ is of dimension $n+2$. It is generated by $n+1$ simple closed loops $\gamma_i= \gamma_i(h)$ which make one turn around the $n+1$ punctures $x_i(h)$ on $\Gamma_h$, $A(x_i(h))-h=0$, as well the loop $\delta(h)$ connecting $0$ and $1$ on the singularized curve $ \Gamma_h^{sing}$. The monodromy of the loop $\delta(h)$ is shown on the Figure~\ref{fig1}.

\begin{figure}[!h]
\includegraphics[scale=0.5]{figures/figure1.pdf} 
%%\includegraphics[width=12cm]{figure1}
\caption{The monodromy of $x_i(h), x_j(h)$ and the loop $\delta(h)$. }\label{fig1}
\end{figure}

As the integral $I(h)$ is constant, it follows that
\[
\int_{\gamma_i(h) - \gamma_j(h)} yq(x) \dx \equiv 0
\]
where $A(x_i(0))=A(x_j(0)) = 0$ and $ x_i(0) = 0$, $x_j(0) = 1$. We have also
\begin{align*}
\int_{\gamma_i(h) - \gamma_j(h)} yq(x) \dx &= \int_{\gamma_i(h) - \gamma_j(h)} \frac{q(x)}{h-A(x)} \dx \\
&= -2\pi i \left(\frac{q(x_i(h))}{A'(x_i(h))} - \frac{q(x_j(h))}{A'(x_j(h))}\right)\\
&= -2\pi i \big[q(x_i(h))x_i'(h) - q(x_j(h)) x_j'(x)\big] \\
&= -2\pi i \frac{\dd}{\ddh} \big[Q(x_i(h)) - Q(x_j(h))\big]\\
\intertext{where}
Q(x) &= \int q(x)\dx
\end{align*}
is a primitive of $q$, and $x_i(h)$ are the roots of the polynomial $A(x)-h$ (we used that $A'(x_i(h)) x_i'(h) \equiv 1$). We denote,
\begin{equation}
J(h) = \int_{x_i(h)-x_j(h)} Q = Q(x_i(h)) - Q(x_j(h))
\end{equation}
and call $J$ an Abelian integral of dimension zero along the zero-cycle
\[
x_i(h)-x_j(h) \in H_0(\{A(x)=h \}, \Z)
\]
(\cite[Definition~1]{gamo07}). If the Abelian integral $I(h)$ vanishes identically, then the same holds true for $J'(h)$, hence $J(h)= \const$ and it is easy to check that the constant is zero, $J(h) \equiv 0$.

The set of rational functions $Q$ such that $ Q(x_i(h)) \equiv Q(x_j(h)$ is a subfield of the field of all rational functions $\Cbb(x)$. By the Lüroth theorem this subfield is of the form $\Cbb(W)$ for suitable rational function $W$. It follows that $Q= \tilde Q \circ W, A = \tilde A \circ W$ where $\tilde Q, \tilde A, W$ are rational functions. As $Q$ is a polynomial, then $Q^{-1}(\infty)= \{ \infty \}$ which implies that $\tilde Q^{-1}(\infty) = \{p\}$ for some $p\in \Pbb^1$ and also $W^{-1}(p) = \{\infty\}$, and similarly $\tilde A^{-1}(\infty) = \{p\}$ Let $\varphi$ be a Mobius functions such that $\varphi(p)=\infty$. Then the functions
\[
\tilde Q \circ \varphi^{-1}, \varphi \circ W, \tilde A \circ \varphi^{-1}
\]
are polynomials. For this reason we may suppose that $\tilde Q, \tilde A, W$ are polynomials. If $W(x_i(h))\equiv W(x_j(h)) $, then clearly $W(0) = W(1)$ and the Theorem~\ref{ccth} is proved. If $W(x_i(h))\not\equiv W(x_j(h)) $, then we denote $ w_i(h) = W(x_i(h)), w_j(h) = W(x_j(h)) $ and note that
\[
\tilde A(w_i(h)) \equiv A(w_j(h)) \equiv h, \tilde Q(w_i(h)) = \tilde Q(w_j(h)).
\]
We may reason then by induction on $\tilde A, \tilde Q$ which have a smaller degree than $A, Q$ respectively. Thus this process must stop and we get $W$ with $W(x_i(h))\equiv W(x_j(h)) $, and hence $W(0)=W(1)$.
\end{proof}



